% night:P1.9
\begin{proof}
\begin{enumerate}[label=\arabic*.]
\item Apply \cite[Theorem~8.7 and Corollary~8.8, pp.~57--59, attached-source edition]{Yang2026}
with $n=3$, $t=q$, $D(q)=C(q)$, $B=I_3$, and the displayed
periodic Kuhn triangulation.  These substitutions give exactly the
action
\[
 (z,q)\longmapsto
 \bigl(q^\lambda\exp(2\pi iC(q)\lambda)z,q\bigr)
\]
and the period order $[I_3\mid sI_3+C(q)]$, where
$q=\exp(2\pi i s)$.  Thus the identity on the dense torus, with
the same integral monomial charts, identifies our quotient with
the cited quotient.  The source theorem supplies its smoothness,
Hausdorffness, properness, and the reduced local equations
$q=x_0\cdots x_k$.  Flatness follows since this nonconstant
function is a nonzero divisor in each smooth chart.
\item For an index-$d$ period sublattice $\Gamma=B\mathbb Z^3$,
the same substitutions give the source's quotient map induced
by $\Gamma\subset\mathbb Z^3$.  It has degree $d$, and
Corollary~8.8 gives $\chi(W)=3!d=6d$ for this marked Kuhn fan.
\end{enumerate}
\end{proof}

% maps:src:F.1.1
\begin{proof}
\begin{enumerate}[label=\arabic*.]
\item The specified rank-three Kuhn fibre is the case $n=3$
of \cite[Corollary~8.8, p.~59, attached-source edition]{Yang2026}: the identical
height-one fan and the inclusion $\Gamma\subset\mathbb Z^3$
identify its index-$d$ quotient, and give central Euler number
$3!d=6d$.
\item The regular fibre has Euler number $\chi(T^6)=0$.
Each other reduced fibre is a free finite quotient of $T^6$,
so its Euler number is also zero.  Euler additivity for the
stated proper triangulable map therefore gives $\chi(M)=6d$.
\end{enumerate}
\end{proof}
